% Test fixture: IEEEtran conference paper
\documentclass[conference]{IEEEtran}
\usepackage{amsmath,amssymb,amsfonts}
\usepackage{algorithmic}
\usepackage{graphicx}
\usepackage{textcomp}
\usepackage{xcolor}
\usepackage{cite}
\usepackage{booktabs}

\begin{document}

\title{A Comprehensive Test Document for WordTex Conversion Fidelity}

\author{
    \IEEEauthorblockN{Jane Doe}
    \IEEEauthorblockA{Department of Computer Science\\
    University of Testing\\
    test@example.edu}
    \and
    \IEEEauthorblockN{John Smith}
    \IEEEauthorblockA{Applied Mathematics Lab\\
    Research Institute\\
    john@example.org}
}

\maketitle

\begin{abstract}
This document serves as a comprehensive test fixture for validating the
bidirectional conversion fidelity of the WordTex system. It exercises all
major LaTeX constructs including mathematical typesetting, tables with
booktabs, figures, algorithms, theorems, and complex cross-referencing.
\end{abstract}

\section{Introduction}
\label{sec:intro}

The problem of converting between \LaTeX{} and Microsoft Word formats
has been a long-standing challenge in academic publishing. This paper
presents a test document that exercises the full range of constructs
that must be preserved during bidirectional conversion.

We refer the reader to Section~\ref{sec:math} for mathematical content
and Table~\ref{tab:results} for experimental results.

\section{Mathematical Content}
\label{sec:math}

\subsection{Inline Mathematics}

The Euler identity $e^{i\pi} + 1 = 0$ combines five fundamental constants.
The quadratic formula gives $x = \frac{-b \pm \sqrt{b^2 - 4ac}}{2a}$.
Consider a function $f: \mathbb{R}^n \to \mathbb{R}$ that is $\mathcal{C}^2$.

\subsection{Display Equations}

The Cauchy-Schwarz inequality:
\begin{equation}
    \label{eq:cauchy-schwarz}
    \left| \sum_{k=1}^{n} a_k b_k \right|^2 \leq
    \left( \sum_{k=1}^{n} a_k^2 \right)
    \left( \sum_{k=1}^{n} b_k^2 \right)
\end{equation}

Maxwell's equations in differential form:
\begin{align}
    \nabla \cdot \mathbf{E} &= \frac{\rho}{\epsilon_0} \label{eq:gauss} \\
    \nabla \cdot \mathbf{B} &= 0 \label{eq:gauss-mag} \\
    \nabla \times \mathbf{E} &= -\frac{\partial \mathbf{B}}{\partial t} \label{eq:faraday} \\
    \nabla \times \mathbf{B} &= \mu_0 \mathbf{J} + \mu_0 \epsilon_0 \frac{\partial \mathbf{E}}{\partial t} \label{eq:ampere}
\end{align}

A piecewise function:
\begin{equation}
    f(x) = \begin{cases}
        x^2 & \text{if } x \geq 0 \\
        -x^2 & \text{if } x < 0
    \end{cases}
\end{equation}

Matrix notation:
\begin{equation}
    \mathbf{A} = \begin{pmatrix}
        a_{11} & a_{12} & \cdots & a_{1n} \\
        a_{21} & a_{22} & \cdots & a_{2n} \\
        \vdots & \vdots & \ddots & \vdots \\
        a_{n1} & a_{n2} & \cdots & a_{nn}
    \end{pmatrix}
\end{equation}

\section{Tables}
\label{sec:tables}

\begin{table}[htbp]
    \caption{Conversion Accuracy Results}
    \label{tab:results}
    \centering
    \begin{tabular}{@{}lccc@{}}
        \toprule
        \textbf{Method} & \textbf{Precision} & \textbf{Recall} & \textbf{F1} \\
        \midrule
        Baseline     & 0.72 & 0.68 & 0.70 \\
        Enhanced     & 0.89 & 0.85 & 0.87 \\
        WordTex SIR  & \textbf{0.98} & \textbf{0.97} & \textbf{0.975} \\
        \bottomrule
    \end{tabular}
\end{table}

\begin{table}[htbp]
    \caption{Multi-column and Multi-row Example}
    \centering
    \begin{tabular}{|l|c|c|c|}
        \hline
        \multirow{2}{*}{\textbf{Category}} & \multicolumn{3}{c|}{\textbf{Metrics}} \\
        \cline{2-4}
        & A & B & C \\
        \hline
        Row 1 & 1.0 & 2.0 & 3.0 \\
        Row 2 & 4.0 & 5.0 & 6.0 \\
        \hline
    \end{tabular}
\end{table}

\section{Lists and Enumerations}

\subsection{Itemized List}
\begin{itemize}
    \item First item with \textbf{bold text}
    \item Second item with \textit{italic text}
    \item Third item with \textsc{Small Caps}
    \item Nested list:
    \begin{itemize}
        \item Sub-item alpha
        \item Sub-item beta
    \end{itemize}
\end{itemize}

\subsection{Enumerated List}
\begin{enumerate}
    \item Parse the input document
    \item Build the Semantic IR
    \item Transform to target format
    \item Validate round-trip fidelity
\end{enumerate}

\section{Code Listing}

\begin{verbatim}
fn convert(doc: &SirDocument) -> Result<Vec<u8>> {
    let pipeline = TransformPipeline::new();
    let ooxml = pipeline.sir_to_ooxml(doc)?;
    Ok(ooxml.to_bytes())
}
\end{verbatim}

\section{Theorem Environments}

\newtheorem{theorem}{Theorem}
\newtheorem{lemma}[theorem]{Lemma}
\newtheorem{definition}{Definition}

\begin{theorem}[Fundamental Theorem of Calculus]
    \label{thm:ftc}
    Let $f$ be a continuous real-valued function on $[a,b]$ and let $F$
    be an antiderivative of $f$. Then
    \begin{equation}
        \int_a^b f(x)\,dx = F(b) - F(a).
    \end{equation}
\end{theorem}

\begin{lemma}
    If $f$ is Riemann integrable on $[a,b]$, then
    $\left|\int_a^b f(x)\,dx\right| \leq \int_a^b |f(x)|\,dx$.
\end{lemma}

\begin{definition}
    A function $f: X \to Y$ is called \emph{continuous} if for every
    open set $U \subseteq Y$, the preimage $f^{-1}(U)$ is open in $X$.
\end{definition}

\section{Cross-References}

As shown in Equation~\eqref{eq:cauchy-schwarz}, and demonstrated
in Theorem~\ref{thm:ftc}, the results in Table~\ref{tab:results}
confirm the effectiveness of the approach described in
Section~\ref{sec:intro}.

\section{Conclusion}

This test document exercises the major LaTeX constructs that the
WordTex conversion system must handle with zero formatting loss.

\bibliographystyle{IEEEtran}
\bibliography{references}

\end{document}
